\documentclass[11pt]{amsart}
\usepackage[T1]{fontenc}
\usepackage{lmodern}
\usepackage{amsmath,amssymb,amsthm}
\usepackage[margin=1.15in]{geometry}
\usepackage[hidelinks]{hyperref}
\hypersetup{pdftitle={Low-intensity Voronoi limits of correlated point processes in the hyperbolic plane},
pdfauthor={Yitzchak Shmalo}}
\newtheorem{theorem}{Theorem}[section]
\newtheorem{lemma}[theorem]{Lemma}
\newtheorem{corollary}[theorem]{Corollary}
\newtheorem{question}[theorem]{Question}
\theoremstyle{definition}
\newtheorem{example}[theorem]{Example}
\theoremstyle{remark}
\newtheorem{remark}[theorem]{Remark}
\newcommand{\HH}{\mathbb{H}}
\newcommand{\D}{\mathbb{D}}
\newcommand{\R}{\mathbb{R}}
\newcommand{\E}{\mathbb{E}}
\newcommand{\Prob}{\mathbb{P}}
\newcommand{\M}{\mathcal{M}}
\newcommand{\Lip}{\mathrm{Lip}}
\newcommand{\Vor}{\mathrm{Vor}}
\newcommand{\area}{A}
\newcommand{\doi}[1]{\href{https://doi.org/#1}{doi:\nolinkurl{#1}}}
\newcommand{\arxiv}[1]{\href{https://arxiv.org/abs/#1}{arXiv:\nolinkurl{#1}}}
\setlength{\emergencystretch}{2.5em}

\title[Low-intensity Voronoi limits of correlated point processes]{Low-intensity Voronoi limits of correlated point
processes in the hyperbolic plane}
\author{Yitzchak Shmalo}
\address{Einstein Institute of Mathematics, The Hebrew University of Jerusalem, Jerusalem, Israel}
\email{yspennstate@gmail.com}
\date{November 2026}
\subjclass[2020]{60G55, 60D05, 60G57, 51M10}
\keywords{ideal Poisson--Voronoi tessellation, hyperbolic plane, low-intensity limit, correlation functions,
determinantal point process, $\alpha$-hyperbolic ensemble}

\raggedbottom
\begin{document}
\begin{abstract}
The Voronoi tessellations of a Poisson process of intensity $\lambda$ on the hyperbolic plane converge, as
$\lambda\to0$, to the ideal Poisson--Voronoi tessellation. D'Achille and Khezeli asked whether other point processes
can produce a different random ideal Voronoi diagram in the same limit, and named the $\alpha$-hyperbolic ensemble, as
$\alpha\to0$, as a natural candidate. We show that this candidate gives the ideal Poisson--Voronoi tessellation again.
More generally, the limit is the ideal Poisson--Voronoi tessellation whenever the $k$-point correlation functions are
bounded by $C^k\lambda^k$ and the normalized correlations $\rho_k/\lambda^k$ tend to one uniformly on configurations whose points are at mutual
distance at least $2\log(1/\lambda)-L$, for every fixed $k$ and $L$. Among families satisfying the bound, a different limit
therefore requires, for some $k$, a relative discrepancy of the $k$-point correlation that does not vanish at distances
$2\log(1/\lambda)+O(1)$.
\end{abstract}
\maketitle

\section{Introduction}

Let $X(\lambda)$ be a Poisson process of intensity $\lambda$ on the hyperbolic plane $\HH^2$. As $\lambda\to0$ the
nuclei of its Voronoi tessellation leave every bounded set, yet the tessellation itself converges. Bhupatiraju~\cite{Bhu19}
proved that the low-intensity limit of Bernoulli--Voronoi measures, with ties broken by labels, on the $d$-regular tree
with $d\ge3$ exists~\cite[Theorem~2.1]{Bhu19}
and studied low-intensity Poisson--Voronoi boundaries in the hyperbolic plane~\cite[Chapter~3]{Bhu19}. Budzinski, Curien
and Petri~\cite{BCP25} studied low-intensity Poisson--Voronoi tessellations of $\HH^2$, used them to bound the
Cheeger constant of closed hyperbolic surfaces of large genus, and conjectured that a limiting random tessellation
exists. D'Achille, Curien, Enriquez, Lyons and \"Unel~\cite{DCE26} proved this and studied the limit, the ideal
Poisson--Voronoi tessellation (IPVT): its cells are indexed by a Poisson process of points of the boundary at
infinity with real delay marks, each point carrying a horofunction shifted by its delay. Fr\k{a}czyk, Mellick and
Wilkens~\cite{FMW23} used low-intensity limits of Poisson--Voronoi tessellations of higher rank symmetric spaces to
prove fixed price one in higher rank. D'Achille and Khezeli~\cite{DK26} gave a necessary and sufficient criterion for
the convergence of the shifted Poisson nuclei on a proper pointed measured metric space, with convergence of the
diagrams under a nondegeneracy condition, and showed that on graphs with counting measure satisfying their growth
conditions Bernoulli and
Poisson nuclei have the same subsequential diagram limits~\cite[Lemma~4.9]{DK26}.

In their list of problems, D'Achille and Khezeli ask~\cite[Question~7.2]{DK26} whether a non-trivial random ideal
Voronoi diagram, different from the IPVT, can be obtained as a low-intensity limit of the Voronoi tessellations of a
point process other than the Poisson process. As a natural candidate in dimension two they propose the
$\alpha$-hyperbolic ensemble, a family of point processes indexed by a parameter $\alpha>0$ studied by Bufetov,
Garc\'{\i}a-Zelada and Lin~\cite{BGZL23}, in the limit $\alpha\to0$. This is the determinantal point process on the
unit disc with the weighted Bergman kernel of weight $(1-|z|^2)^{\alpha-1}$. These invariant determinantal processes
were studied by Jancovici and T\'ellez~\cite{JT98} on the hyperbolic plane; the kernel for every real $\alpha>0$ is
in~\cite[Section~4.3.11]{HKPV09}, and for integer $\alpha$ the process is the set of singular points of a
matrix-valued Gaussian analytic function, by the theorem of Krishnapur~\cite{Kri09}. For $\alpha=1$ it is the zero
set of the Gaussian power series with independent standard complex coefficients, by the theorem of Peres and
Vir\'ag~\cite{PV05}. Its intensity is $\alpha/(4\pi)$ with respect to hyperbolic area, so $\alpha\to0$ is a
low-intensity limit.

We show that the candidate gives the IPVT again. This settles the proposed candidate; whether some other point process
gives a different limit remains open. The reason is general. Write $F(r)=2\pi(\cosh r-1)$ for the area of
a hyperbolic disc of radius $r$, let $o$ be a fixed point, and for small $\lambda>0$ let $t=t(\lambda)$ be the radius
with $\lambda F(t)=1$, so that $t=\log(1/\lambda)+O(1)$.

\begin{theorem}\label{thm:main}
For $0<\lambda<\lambda_0$ let $\Phi_\lambda$ be a simple point process on $\HH^2$ with correlation functions
$\rho_k=\rho^\lambda_k$ with respect to hyperbolic area, with $\rho_1\equiv\lambda$. Assume:
\begin{enumerate}
\item[(A)] there is a constant $C\ge1$ with $\rho_k\le C^k\lambda^k$ for all $k$ and $\lambda$;
\item[(B)] for every $k\ge2$, $b\in\R$ and $L>0$,
\[
  \sup\Bigl\{\Bigl|\frac{\rho_k(x_1,\dots,x_k)}{\lambda^k}-1\Bigr| : x_i\in B(o,t+b),\ d(x_i,x_j)\ge 2t-L
  \text{ for }i\ne j\Bigr\}\longrightarrow0\qquad(\lambda\to0).
\]
\end{enumerate}
Then, as $\lambda\to0$, the Voronoi tessellation of $\Phi_\lambda$ converges in distribution to the IPVT of $\HH^2$, and
the shifted nuclei of $\Phi_\lambda$ converge to the same Poisson process as those of the Poisson process of intensity
$\lambda$.
\end{theorem}

The precise meaning of the convergence is given in Section~\ref{sec:setting}. The bound (A) controls the
configurations excluded in (B) and gives the exponential moments used in Step~2 of the proof. The theorem does not
assume that $\Phi_\lambda$ is invariant under the
isometries of $\HH^2$.

\begin{corollary}\label{cor:ensemble}
As $\alpha\to0$, the Voronoi tessellation of the $\alpha$-hyperbolic ensemble converges in distribution to the IPVT of
$\HH^2$.
\end{corollary}

The nuclei that shape the cells near $o$ lie at distance $t+O(1)$ from it, and two of them are typically at distance
$2t+O(1)$ from each other. Theorem~\ref{thm:main} therefore says that, among families satisfying (A), a different
limit needs, for some $k$, normalized $k$-point correlations that stay away from one at that scale. In the
$\alpha$-hyperbolic ensemble the normalized correlations approach one exponentially fast in the distance, at a rate
that does not depend on $\lambda$. Section~\ref{sec:needs} discusses the constructions that fail and the hard-core case
that remains open.

The proof works on the space of shifted nuclei introduced in~\cite{DCE26,FMW23} with the topology of
D'Achille and Khezeli~\cite[Section~3.1]{DK26}. We show that the shifted nuclei processes converge to the same Poisson
limit as in the Poisson case, by the method of factorial moments; the passage from nuclei to tessellations is then
the argument of~\cite[Proposition~3.3]{DK26}, which does not use the Poisson property of the approximating processes.

\section{The setting}\label{sec:setting}

We use the Poincar\'e disc $\D$ with the metric $4|dz|^2/(1-|z|^2)^2$ of curvature $-1$, hyperbolic area $\area$, and
origin $o=0$. Balls are closed, $B(x,r)$, and
\[
  F(r)=\area(B(o,r))=2\pi(\cosh r-1).
\]
For $0<\lambda<1/F(2)$ let $t=t(\lambda)$ be the radius with $\lambda F(t)=1$, the shift of~\cite[(3.1)]{DK26}.
Since $F(r)\ge \pi e^r/2$ for $r\ge2$, we have $t\ge2$ and
\begin{equation}\label{eq:t}
  \lambda e^{t}\le 2/\pi,\qquad t=\log(1/\lambda)-\log\pi+o(1)\quad(\lambda\to0).
\end{equation}

\emph{Shifted nuclei.} Let $\Lip^1$ be the set of $1$-Lipschitz functions on $\D$ and $\Lip^1_0$ the subset of those
vanishing at $o$, with the topology of uniform convergence on compact sets; $\Lip^1_0$ is compact. A function
$f\in\Lip^1$ is identified with the pair $(f-f(o),f(o))\in\Lip^1_0\times\R$, and $f(o)$ is called its delay. Put
$S=\Lip^1_0\times[-\infty,\infty)$ and let $\M$ be the space of Radon measures on $S$ with the vague topology. A set is
relatively compact in $S$ exactly when it is contained in $\Lip^1_0\times[-\infty,b]$ for some $b\in\R$. For $s\in\R$
and $x\in\D$ let $\iota_s(x)=d(x,\cdot)-s\in\Lip^1$, with delay $d(o,x)-s$; for a locally finite set $\varphi\subset\D$
let $\iota_s\varphi=\sum_{x\in\varphi}\delta_{\iota_s(x)}$. Since
\begin{equation}\label{eq:preimage}
  \iota_s^{-1}\bigl(\Lip^1_0\times[-\infty,b]\bigr)=B(o,s+b),
\end{equation}
with $B(o,r)=\emptyset$ for $r<0$, $\iota_s\varphi$ belongs to $\M$.

\emph{Voronoi diagrams.} For a locally finite family $(f_i)$ in $\Lip^1$ with delays tending to $+\infty$ the cells are
$C_i=\{y\in\D: f_i(y)\le f_j(y)\text{ for all }j\}$, and $\Vor((f_i))$ is the collection of the nonempty cells (for a
general family some cells can be empty, and these are discarded), a point process on the space $\mathcal F$ of
nonempty closed subsets of $\D$ with the Fell topology~\cite[Section~2.2]{DK26}; point measures on $\mathcal F$ carry
the vague topology. To record which function generates which cell we also use the marked process
$\sum_i\delta_{(f_i,C_i)}$ on $S\times\mathcal F$, over the $i$ with $C_i\ne\emptyset$, with the vague topology on the
product. Subtracting a constant from every $f_i$ does not change the cells, so $\Vor(\iota_s\varphi)=
\Vor(\varphi)$ for every $s$, where $\Vor(\varphi)$ is the usual Voronoi tessellation of $\varphi$.

\emph{The Poisson limit.} For $\theta\in\partial\D$ let $\beta_\theta$ be the Busemann function of $\theta$ normalized by
$\beta_\theta(o)=0$. Let $\mu$ be the image of the measure $(2\pi)^{-1}d\theta\otimes e^{s}ds$ on $\partial\D\times\R$
under $(\theta,s)\mapsto(\beta_\theta,s)$, and let $\Pi$ be a Poisson process on $S$ with intensity $\mu$. Put
$\mu_\lambda=\lambda\,(\iota_{t(\lambda)})_\#\area$, the intensity measure of $\iota_{t(\lambda)}X(\lambda)$. We use three
facts.

\begin{enumerate}
\item[(F1)] $\mu_\lambda\to\mu$ in $\M$ as $\lambda\to0$, and $\mu$ is diffuse. This is the full convergence
of~\cite[Theorem~1.2]{DK26} on $\HH^2$, see the proof of~\cite[Proposition~1.1]{DK26} and~\cite{DCE26}. As a check,
$\mu_\lambda(\Lip^1_0\times[-\infty,b])=\lambda F(t+b)\to e^b=\mu(\Lip^1_0\times[-\infty,b])$.
\item[(F2)] Almost surely the delays of the points of $\Pi$ are distinct and $\Vor(\Pi)$ is nondegenerate, that is,
every cell is the closure of the set where the defining inequalities are strict. The delays are distinct because the
image of $\mu$ on the delay coordinate is $e^s\,ds$, which has no atoms. Nondegeneracy is~\cite[Corollary~5.4]{DK26}
together with~\cite[Lemma~5.3]{DK26}, since $\HH^2$ is a CAT(0) Riemannian manifold and $\mu$ is diffuse. The random
diagram $\Vor(\Pi)$ is the IPVT of $\HH^2$, and $\Vor(X(\lambda))\to\Vor(\Pi)$~\cite{DCE26},~\cite[Proposition~1.1]{DK26}.
\item[(F3)] The passage from nuclei to diagrams, Lemma~\ref{lem:transfer} below.
\end{enumerate}

\begin{lemma}\label{lem:transfer}
Let $t_n\to\infty$ and let $\Phi_n$ be random locally finite subsets of $\D$ such that $\iota_{t_n}\Phi_n\to\Pi$ in
distribution in $\M$. Then $\Vor(\Phi_n)\to\Vor(\Pi)$ in distribution as point processes on $\mathcal F$, and the
marked processes $\sum_{x\in\Phi_n}\delta_{(\iota_{t_n}(x),C_x)}$ converge in distribution to
$\sum_i\delta_{(f_i,C_i)}$ as point processes on $S\times\mathcal F$, where $(f_i)$ are the points of $\Pi$, $C_x$ and
$C_i$ are the corresponding cells, and empty cells are discarded in both.
\end{lemma}

\begin{proof}
This is part~(i) of~\cite[Proposition~3.3]{DK26} with $X(\lambda_n)$ replaced by $\Phi_n$. That proof applies
Skorokhod's representation theorem to the convergence of the shifted nuclei processes in $\M$ and then argues on a
fixed outcome: for each $N$ and all large $n$ the first $N$ points of $\Pi$, ordered by delay, are limits of points of
$\iota_{t_n}\Phi_n$, and a slight modification of~\cite[Theorem~2.3]{DCE26} gives the convergence of the diagrams on the
event of (F2). The Poisson property of $X(\lambda_n)$ is used only to obtain the convergence in $\M$, through part~(iv)
of~\cite[Proposition~3.2]{DK26}, and here that convergence is the hypothesis. The limit is the same process $\Pi$, for
which (F2) holds.

Only finitely many functions matter on a bounded set, and this makes the argument a finite one. Let $f_1$ have the
smallest delay $s_1$ of a family. Since the functions are $1$-Lipschitz, on $B(o,r)$ every $f_i$ with delay $s_i>s_1+2r$
satisfies $f_i\ge s_i-r>s_1+r\ge f_1$, so its cell does not meet $B(o,r)$. On the coupled outcome choose $c>s_1+2r+1$
with no point of $\Pi$ at delay $c$. The sets $\Lip^1_0\times[-\infty,c]$ are compact in $S$, so for large $n$ the
points of $\iota_{t_n}\Phi_n$ with delay at most $c$ are finitely many, as many as those of $\Pi$, and converge to them
one by one; in particular the smallest delays converge, and no point of small delay can appear unseen. Only these
points determine the cells on $B(o,r)$. Strict inequalities between finitely many of them hold with a positive margin
on compact subsets of the set where they are strict and survive small perturbations, and weak inequalities pass to the limit; by nondegeneracy this gives
the local Fell convergence of the cells on $B(o,r)$, and the convergence of the marked processes.
The marks used in~\cite{DK26} to break ties between equal delays play no role, since the delays of $\Pi$ are distinct.
\end{proof}

\emph{Correlation functions.} A simple point process $\Phi$ on $\D$ has correlation functions $\rho_k$ with respect to
$\area$ if for all disjoint Borel sets $P_1,\dots,P_m$ and integers $k_1,\dots,k_m\ge0$ with $K=k_1+\dots+k_m$,
\begin{equation}\label{eq:corr}
  \E\prod_{i=1}^m\bigl(\Phi(P_i)\bigr)_{k_i}=\int_{P_1^{k_1}\times\dots\times P_m^{k_m}}\rho_K\,d\area^{K},
\end{equation}
where $(n)_k=n(n-1)\cdots(n-k+1)$, with $(n)_0=1$ and $\rho_0=1$; see~\cite[Section~1.2]{HKPV09}. The Poisson process of intensity $\lambda$ has
$\rho_k\equiv\lambda^k$.

\section{Proof of Theorem~\ref{thm:main}}\label{sec:main}

In the notation of Section~\ref{sec:setting}, we prove that $\iota_{t(\lambda)}\Phi_\lambda\to\Pi$ in distribution in
$\M$ as $\lambda\to0$. Lemma~\ref{lem:transfer} then gives the convergence of $\Vor(\Phi_\lambda)$ to the IPVT, jointly
with the shifted nuclei. Condition (B) refers to the origin; for invariant processes the choice of the origin does not
matter. The first step is a bound on the pairs of points of a large ball that are closer than $2t-L$.

\begin{lemma}\label{lem:pairs}
For $T\ge1$ and $R>0$,
\[
  m(T,R):=\area\otimes\area\bigl\{(x,y)\in B(o,T)^2: d(x,y)<R\bigr\}\le 2F(1)F(T)+12\pi^2e^{T+R/2}.
\]
\end{lemma}

\begin{proof}
The pairs in which one point lies in $B(o,1)$ have measure at most $2F(1)F(T)$. Write the other points in polar
coordinates, $x=(r,\omega)$ and $y=(r',\omega')$ with $1\le r,r'\le T$, where the area element is
$\sinh r\,dr\,d\omega$, and
let $\psi\in(-\pi,\pi]$ be the angle between them. By the hyperbolic law of cosines,
\[
  e^{d(x,y)}\ge\cosh d(x,y)=\cosh r\cosh r'-\sinh r\sinh r'\cos\psi\ge 2\sinh r\sinh r'\sin^2(\psi/2).
\]
For $r\ge1$ we have $\sinh r\ge e^r/3$, and $|\sin(\psi/2)|\ge|\psi|/\pi$, so $d(x,y)<R$ forces
$\psi^2<\tfrac92\pi^2e^{R-r-r'}$ and hence $|\psi|<3\pi e^{(R-r-r')/2}$. For fixed $r,r',\omega$ the set of such $\omega'$
has measure at most $6\pi e^{(R-r-r')/2}$. Using $\sinh r\le e^r/2$,
\[
  \int_1^T\!\!\int_1^T 2\pi\cdot6\pi e^{(R-r-r')/2}\sinh r\sinh r'\,dr\,dr'\le 3\pi^2e^{R/2}\Bigl(\int_0^Te^{r/2}dr\Bigr)^2
  \le12\pi^2e^{T+R/2}.\qedhere
\]
\end{proof}

\begin{proof}[Proof of Theorem~\ref{thm:main}]
Write $\xi_\lambda=\iota_{t}\Phi_\lambda$. Fix $b\in\R$ and disjoint Borel sets $A_1,\dots,A_m\subset
\Lip^1_0\times[-\infty,b]$ with $\mu(\partial A_i)=0$. Let $P_i=\iota_t^{-1}(A_i)$; by~\eqref{eq:preimage} these are
disjoint subsets of $B=B(o,t+b)$, and $\xi_\lambda(A_i)=\Phi_\lambda(P_i)$, $\mu_\lambda(A_i)=\lambda\area(P_i)$.

\emph{Step 1: factorial moments.} Fix $k_1,\dots,k_m$ with $K=\sum k_i\ge1$, let $Q=P_1^{k_1}\times\dots\times
P_m^{k_m}\subset B^K$, and for $L>0$ let $G_L\subset Q$ be the set of $K$-tuples whose points are pairwise at distance at
least $2t-L$. By~\eqref{eq:corr},
\[
  \Delta_\lambda:=\E\prod_i\bigl(\xi_\lambda(A_i)\bigr)_{k_i}-\prod_i\mu_\lambda(A_i)^{k_i}=\int_Q(\rho_K-\lambda^K)\,d\area^K.
\]
For $K=1$ this vanishes because $\rho_1\equiv\lambda$. For $K\ge2$ we split $Q$ into $G_L$ and $Q\setminus G_L$. On $G_L$,
condition (B) gives
\[
  \Bigl|\int_{G_L}(\rho_K-\lambda^K)\Bigr|\le\varepsilon_\lambda\,(\lambda F(t+b))^K,\qquad\varepsilon_\lambda\to0.
\]
A tuple of $Q\setminus G_L$ has a pair of coordinates at distance less than $2t-L$, so by Lemma~\ref{lem:pairs},
\begin{align*}
  \area^K(Q\setminus G_L)&\le\binom K2F(t+b)^{K-2}\,m(t+b,2t-L)\\
  &\le\binom K2F(t+b)^{K-2}\bigl(2F(1)F(t+b)+12\pi^2e^{2t+b-L/2}\bigr).
\end{align*}
By (A) the integrand is at most $(C^K+1)\lambda^K$ in absolute value there. With $\lambda F(t+b)=F(t+b)/F(t)\to e^b$,
$\lambda e^t\le2/\pi$ from~\eqref{eq:t}, and $\lambda F(1)\to0$, we get
\[
  \limsup_{\lambda\to0}|\Delta_\lambda|\le(C^K+1)\binom K2e^{(K-2)b}\cdot48\,e^{b-L/2}.
\]
Since $L$ is arbitrary, $\Delta_\lambda\to0$. By (F1), $\mu_\lambda(A_i)\to\mu(A_i)$, so
\begin{equation}\label{eq:fm}
  \E\prod_i\bigl(\xi_\lambda(A_i)\bigr)_{k_i}\longrightarrow\prod_i\mu(A_i)^{k_i}.
\end{equation}

\emph{Step 2: the counts converge.} By~\eqref{eq:corr} and (A), the left side of~\eqref{eq:fm} is at most
$(C\lambda F(t+b))^K\le c^K$ with $c=Ce^{b+1}$, for all small $\lambda$. For $N=(N_1,\dots,N_m)$ with values in
$\{0,1,\dots\}^m$ and $z\in[0,1]^m$,
\[
  \E\prod_iz_i^{N_i}=\sum_{k_1,\dots,k_m\ge0}\E\prod_i(N_i)_{k_i}\prod_i\frac{(z_i-1)^{k_i}}{k_i!},
\]
which follows from $z^n=\sum_k(n)_k(z-1)^k/k!$ by Fubini's theorem, because
\[
  \sum_{k_1,\dots,k_m\ge0}\E\prod_i(N_i)_{k_i}\prod_i\frac1{k_i!}=\E\,2^{N_1+\dots+N_m},
\]
which is finite when the factorial moments are bounded by $c^K$. Apply this to $N=(\xi_\lambda(A_1),\dots,\xi_\lambda(A_m))$. The terms are bounded by $c^K/\prod_ik_i!$, which is
summable, so by~\eqref{eq:fm} and dominated convergence the generating functions converge on $[0,1]^m$ to
$\prod_i\exp(\mu(A_i)(z_i-1))$. Hence $(\xi_\lambda(A_1),\dots,\xi_\lambda(A_m))$ converges in distribution to
independent Poisson variables with means $\mu(A_i)$, which is the law of $(\Pi(A_1),\dots,\Pi(A_m))$.

\emph{Step 3: convergence in $\M$.} Let $\mathcal B$ be the class of relatively compact Borel sets $A\subset S$ with
$\mu(\partial A)=0$, equivalently $\Pi(\partial A)=0$ almost surely. Any finitely many sets of $\mathcal B$ are finite
unions of disjoint sets of $\mathcal B$ (the atoms of the algebra they generate, whose boundaries lie in the union of
the boundaries), so Step 2 gives the joint convergence in distribution of $(\xi_\lambda(A_1),\dots,\xi_\lambda(A_m))$ for
all $A_1,\dots,A_m\in\mathcal B$. Since $S$ is locally compact, second countable and Hausdorff, this implies
$\xi_\lambda\to\Pi$ in distribution in $\M$~\cite[Theorem~4.11]{Kal17}. The convergence of the diagrams follows from
Lemma~\ref{lem:transfer}, applied along any sequence $\lambda_n\to0$.
\end{proof}

\begin{remark}
Because the sets $\Lip^1_0\times[-\infty,b]$ are compact in $S$, the convergence includes the points of small delay:
in the limit no nucleus comes close to $o$. This is the part of the topology of~\cite[Section~3.1]{DK26} that is
vague at $+\infty$ and weak at $-\infty$, and it is what Lemma~\ref{lem:transfer} needs.
\end{remark}

\section{\texorpdfstring{The $\alpha$-hyperbolic ensemble}{The alpha-hyperbolic ensemble}}\label{sec:ensemble}

Let $m$ be Lebesgue measure on $\D$ and, for $\alpha>0$, $q_\alpha(z)=\frac\alpha\pi(1-|z|^2)^{\alpha-1}$. The
orthogonal projection of $L^2(\D,q_\alpha\,dm)$ onto its subspace of holomorphic functions has kernel
$(1-z\bar w)^{-\alpha-1}$~\cite[Section~4.3.11]{HKPV09}, and the $\alpha$-hyperbolic ensemble $\mathcal H_\alpha$ is the
corresponding determinantal point process on $\D$~\cite[Definition~1.2]{BGZL23}. Written with respect to $m$, its
kernel is
\[
  K_\alpha(z,w)=\sqrt{q_\alpha(z)}\,\frac{1}{(1-z\bar w)^{\alpha+1}}\,\sqrt{q_\alpha(w)}
  =\frac\alpha\pi\,\frac{(1-|z|^2)^{(\alpha-1)/2}(1-|w|^2)^{(\alpha-1)/2}}{(1-z\bar w)^{\alpha+1}},
\]
which is Hermitian and positive semidefinite, and the process is simple~\cite[Chapter~4]{HKPV09}. We prove Corollary~\ref{cor:ensemble} in the stronger form
$\iota_{t(\alpha/4\pi)}\mathcal H_\alpha\to\Pi$ in distribution in $\M$.

\begin{proof}[Proof of Corollary~\ref{cor:ensemble}]
With respect to hyperbolic area $d\area=4\,dm/(1-|z|^2)^2$ the kernel is
$K^\area_\alpha(z,w)=\tfrac14(1-|z|^2)(1-|w|^2)K_\alpha(z,w)$, and its diagonal is the constant $\lambda=\alpha/(4\pi)$.
The correlation functions are $\rho_k(x_1,\dots,x_k)=\lambda^k\det M$ with $M_{ij}=K^\area_\alpha(x_i,x_j)/\lambda$. The
matrix $M$ is positive semidefinite with unit diagonal, and from
$|1-z\bar w|^2-|z-w|^2=(1-|z|^2)(1-|w|^2)$ and $|z-w|/|1-z\bar w|=\tanh(d(z,w)/2)$,
\[
  |M_{ij}|^2=\Bigl(\frac{(1-|x_i|^2)(1-|x_j|^2)}{|1-x_i\bar x_j|^2}\Bigr)^{\alpha+1}
  =\operatorname{sech}^{2(\alpha+1)}\bigl(d(x_i,x_j)/2\bigr).
\]
By Hadamard's inequality~\cite[Theorem~7.8.1, p.~505]{HJ13}, $\det M\le\prod_iM_{ii}=1$, so (A) holds with $C=1$. Every
permutation other than the identity moves at least two indices, and $|M_{ij}|\le1$, so
\[
  |\det M-1|\le\sum_{\sigma\ne\mathrm{id}}\prod_i|M_{i\sigma(i)}|\le k!\max_{i\ne j}|M_{ij}|^2.
\]
If $d(x_i,x_j)\ge2t-L$, then $|M_{ij}|\le\operatorname{sech}^{\alpha+1}(t-L/2)\le2^{\alpha+1}e^{-(\alpha+1)(t-L/2)}$, which
tends to $0$ as $\lambda\to0$, uniformly in the configuration. This is (B). Since $\lambda\to0$ exactly when $\alpha\to0$,
Theorem~\ref{thm:main} applies.
\end{proof}

\begin{remark}
Bufetov, Garc\'{\i}a-Zelada and Lin study $\mathcal H_\alpha$ for fixed $\alpha$ on discs of radius tending to one; the
limit $\alpha\to0$ considered here is a different regime. The process is invariant under the automorphisms of $\D$,
since $|M_{ij}|$ depends only on the distances and an automorphism changes $K^\area_\alpha$ by a factor
$c(z)\bar c(w)$ with $|c|=1$, which leaves determinants unchanged. The theorem does not need this.
\end{remark}

\section{What a different limit requires}\label{sec:needs}

Let $x,y\in\D$ lie at distances $t+s$ and $t+s'$ from $o$, at angle $\psi$. As $t\to\infty$ with $s,s'\in\R$ and
$\psi\not\equiv0\pmod{2\pi}$ fixed, the law of cosines gives
\begin{equation}\label{eq:limitdist}
  d(x,y)=2t+s+s'+2\log|\sin(\psi/2)|+o(1).
\end{equation}
The expansion is not uniform as $\psi\to0$; the proof does not need it there, since Lemma~\ref{lem:pairs} handles close
pairs. By~\eqref{eq:t}, $2t=2\log(1/\lambda)+O(1)$. The nuclei that matter for the cells meeting a fixed ball have delays
$s=O(1)$, so their mutual distances are $2\log(1/\lambda)+O(1)$. Theorem~\ref{thm:main} therefore reads as follows.
Suppose a family satisfies (A) and its Voronoi tessellations have a limit other than the IPVT, or no limit. Then for
some $k$, $b$ and $L$, along a sequence $\lambda_n\to0$, the supremum in (B) stays bounded away from zero: the relative
error of the $k$-point correlation does not vanish on configurations in $B(o,t+b)$ with mutual distances at least
$2t-L$. This need not be the two-point correlation. Under (A), the IPVT is the limit whenever, for every $k$,
$|\rho_k/\lambda^k-1|\le h_k(\min_{i\ne j}d(x_i,x_j))$ with $h_k$ independent of $\lambda$ and $h_k(r)\to0$ as
$r\to\infty$, and also whenever the relative errors tend to zero uniformly beyond distance $(2-\delta)\log(1/\lambda)$
for some $\delta>0$, since $2t-L$ eventually exceeds that distance. A nonincreasing envelope $h_{k,\lambda}$ that depends on
$\lambda$ also suffices if $h_{k,\lambda}(2t-L)\to0$ for every fixed $L>0$. Decay of absolute errors, or relative decay for each fixed intensity
without uniform control at this scale, does not by itself imply (B).

The condition is necessary for a different limit, not sufficient, as the simplest example shows.

\begin{example}[Mixed Poisson processes]
Let $W$ be a random variable with values in $[w_0,w_1]$, $0<w_0<w_1$, with $\E W=1$, that is not constant, and given
$W$ let $\Phi_\lambda$ be a Poisson process of intensity $\lambda W$. Then $\rho_k=\lambda^k\E W^k$, so (A) holds with
$C=w_1$ and (B) fails for $k=2$, since $\E W^2>1$. Given $W=w$, the process $\iota_{t-\log w}\Phi_\lambda$ is Poisson with
intensity measure $w\,\lambda(\iota_{t-\log w})_\#\area$, the image of $w\mu_\lambda$ under the shift of every delay by
$\log w$, which converges to $\mu$ by (F1), since the factor $e^{-\log w}$ from translating the density $e^s$ cancels
$w$; so $\iota_{t-\log w}\Phi_\lambda\to\Pi$ by Steps~2 and~3 of the proof of Theorem~\ref{thm:main}, applied to the
Poisson case. Lemma~\ref{lem:transfer} gives $\Vor(\Phi_\lambda)\to\Vor(\Pi)$
given $W=w$, for every $w$, and hence unconditionally. So the limit is the IPVT.
\end{example}

Short-range clustering does not change the limit either.

\begin{example}[Poisson cluster processes]
Take Poisson parents of intensity $\kappa\lambda$, with $\kappa>0$ fixed, and attach to them independent identically
distributed decorations, independent of the parents, whose fixed law is isometry-covariant; each decoration is a
nonempty finite simple configuration within a fixed distance $\Delta$ of its parent, with finite mean size $\E N$.
The offspring then have intensity $\kappa\lambda\E N$, and $\kappa=(\E N)^{-1}$ keeps it equal to $\lambda$. An offspring at distance
$\delta\le\Delta$ from its parent, at angle $\omega$ to the parent's outward radial direction, has limiting delay
$s+V$ with $V=\log(\cosh\delta+\sinh\delta\cos\omega)$ by the law of cosines, so a parent with limiting coordinates
$(\theta,s)$ has offspring at $(\theta,s+V_j)$ with $|V_j|\le\Delta$, and the law of $V=\min_jV_j$ does not depend on
$(\theta,s)$. Only parents with $s\le b+\Delta$ have offspring of delay at most $b$, so the limit of the shifted
nuclei is a locally finite Poisson cluster process. Offspring of one parent share the direction $\theta$, so their limiting functions differ by constants
and only the offspring of least delay has a nonempty cell (ties have probability zero); the other offspring are
eventually dominated on every fixed ball, so the finite comparison in the proof of Lemma~\ref{lem:transfer} carries the
diagram limit from the representatives to all the offspring. By the marked mapping theorem
these representatives form a Poisson process of intensity $a\,e^s\,ds\,d\theta/(2\pi)$, where
\[
  \kappa\int e^{s}\,\Prob\{s+V\in dv\}\,ds=\kappa\,\E[e^{-V}]\,e^{v}\,dv,\qquad a=\kappa\,\E[e^{-V}]\in(0,\infty),
\]
and shifting every delay by $\log a$ gives intensity $\mu$ and preserves every cell, so the diagram is again the IPVT.
\end{example}

The natural ways of producing long-range dependence lose either the dependence or the intensity at the scale
$2t$. For instance, the Mat\'ern hard-core process of type II with exclusion distance $R$, obtained from a Poisson process
of any intensity $\nu$ with independent uniform marks by keeping the points with no point of smaller mark within
distance $R$, has intensity $\nu\int_0^1e^{-\nu F(R)u}\,du=(1-e^{-\nu F(R)})/F(R)\le1/F(R)$. For $R=2t-L$ this is at most
of order $\lambda^2e^{L}$, far below $\lambda$. The
following case is not covered by the theorem.

\begin{question}\label{q:hardcore}
By~\eqref{eq:limitdist}, an exclusion distance $2t-L$ becomes, in the limit coordinates $(\theta,s)\in\partial\D\times\R$,
the condition
\[
  s_i+s_j+2\log|\sin((\theta_i-\theta_j)/2)|\ge-L\qquad(i\ne j)
\]
on every pair of nuclei, which says that the horoballs $\{x:\beta_{\theta_i}(x)+s_i\le-L/2\}$, the limits of the
balls of radius $t-L/2$ about the original nuclei, have pairwise disjoint interiors. Is there a family of isometry-invariant point processes of intensity of
order $\lambda$ on $\HH^2$, with exclusion distance $2t-L$, whose shifted nuclei converge to a hard-core process on
$\partial\D\times\R$ with this constraint, and is the Voronoi diagram of such a limit different from the IPVT?
\end{question}

Under $\Pi$ the expected number of pairs of points with delays at most $b$ that violate the constraint is finite, by
the same computation as in the proof of Lemma~\ref{lem:pairs}, and it is positive. A hard-core limit of this kind
therefore differs from $\Pi$ in law on every window of delays $[-\infty,b]$; whether its Voronoi diagram differs from the
IPVT is the content of Question~\ref{q:hardcore}.

\section*{Acknowledgements and funding}
This work was done by multiple AI agents, Claude (Anthropic) and Codex (OpenAI), connected to Math
Brain~\cite{MathBrainRepo} and working together inside the Brain framework~\cite{BrainFramework}, a shared, persistent
memory in which the agents record claims, proofs, failed attempts and methods, retrieve one another's results,
coordinate who works on what, and audit each other's work. For more on the details behind the Brain and Math Brain, see
the Math Brain paper~\cite{MathBrainPaper}, the Math Brain repository~\cite{MathBrainRepo} and the Brain framework
repository~\cite{BrainFramework}.

The research presented in this paper was supported by the European Research Council (ERC) under the European Union's
Horizon Europe research and innovation program (grant agreement No.~101041711), by the Simons Foundation as part of
the Collaboration on the Mathematical and Scientific Foundations of Deep Learning, by Heights Labs, by Convex Nexus
Capital, by the Israel Science Foundation (grant number 2258/19), and by the Israel Science Foundation (ISF Grant
4101/25). Computational resources were provided in part by Google Cloud through its education credits program (credit
530550103).

\section*{Data availability}
No data were generated or analysed for this paper. The \LaTeX{} source of this version is available at\\
\url{https://github.com/yspennstate/hyperbolic-voronoi-correlated-nuclei/tree/v1.1.1}.

The paper is part of the Math Brain project, whose public repository is given in~\cite{MathBrainRepo}.

\section*{Declaration of competing interest}
The author holds equity in Heights Labs and in Convex Nexus Capital and has a management role at Convex Nexus Capital.
These entities are included in the funding declaration above.

\section*{Declaration of generative AI and AI-assisted technologies}
Claude agents (Anthropic) proved Theorem~\ref{thm:main}, Lemma~\ref{lem:pairs} and Corollary~\ref{cor:ensemble} and
prepared the manuscript; Codex agents (OpenAI) refereed the proofs. The question and the candidate process are due to
Matteo D'Achille and Ali Khezeli~\cite{DK26}, whose framework is used throughout. The ideal Poisson--Voronoi
tessellation of $\HH^2$ was conjectured by Thomas Budzinski, Nicolas Curien and Bram Petri~\cite{BCP25} and constructed
by Matteo D'Achille, Nicolas Curien, Nathana\"el Enriquez, Russell Lyons and Meltem \"Unel~\cite{DCE26}; Sandeep
Bhupatiraju~\cite{Bhu19} studied the low-intensity limit earlier on regular trees and in the hyperbolic plane. The
$\alpha$-hyperbolic ensemble is the one studied by Alexander Bufetov, David Garc\'{\i}a-Zelada and Zhaofeng
Lin~\cite{BGZL23}; these determinantal processes were studied by Jancovici and T\'ellez~\cite{JT98} and, for integer
$\alpha$, through matrix Gaussian analytic functions, by Manjunath Krishnapur~\cite{Kri09}. New here are the correlation
criterion of Theorem~\ref{thm:main} and its application in Corollary~\ref{cor:ensemble}. The author directed the work
and is responsible for the content of the paper and for decisions concerning its dissemination.

\begin{thebibliography}{99}
\bibitem{BGZL23} A.~I.~Bufetov, D.~Garc\'{\i}a-Zelada and Z.~Lin, \emph{Fluctuations of the process of moduli for the
Ginibre and hyperbolic ensembles}, J. Math. Phys. Anal. Geom. \textbf{19} (2023), no.~1, 74--106,
\doi{10.15407/mag19.01.074}, \arxiv{2202.11687}.
\bibitem{Bhu19} S.~Bhupatiraju, \emph{The low-intensity limit of Bernoulli--Voronoi and Poisson--Voronoi measures},
Ph.D. thesis, Indiana University, Bloomington, 2019,
\url{https://scholarworks.iu.edu/dspace/items/0af0462b-d7ff-4ef4-98b1-7f93bb59fe32}.
\bibitem{BCP25} T.~Budzinski, N.~Curien and B.~Petri, \emph{On Cheeger constants of hyperbolic surfaces}, Invent.
Math. \textbf{242} (2025), no.~2, 511--530, \doi{10.1007/s00222-025-01361-w}, \arxiv{2207.00469}.
\bibitem{DCE26} M.~D'Achille, N.~Curien, N.~Enriquez, R.~Lyons and M.~\"Unel, \emph{Ideal Poisson--Voronoi
tessellations on hyperbolic spaces}, Ann. Probab. \textbf{54} (2026), no.~2, 846--891, \doi{10.1214/25-AOP1781},
\arxiv{2303.16831}.
\bibitem{DK26} M.~D'Achille and A.~Khezeli, \emph{Convergence towards ideal Poisson--Voronoi tessellations with a focus
on Diestel--Leader graphs}, preprint (2026), \arxiv{2606.30517v2}.
\bibitem{FMW23} M.~Fr\k{a}czyk, S.~Mellick and A.~Wilkens, \emph{Poisson--Voronoi tessellations and fixed price in higher
rank}, Ann. of Math., to appear; preprint (2023), \arxiv{2307.01194}.
\bibitem{HJ13} R.~A.~Horn and C.~R.~Johnson, \emph{Matrix Analysis}, 2nd ed., Cambridge University Press, Cambridge,
2013, \doi{10.1017/CBO9781139020411}.
\bibitem{HKPV09} J.~B.~Hough, M.~Krishnapur, Y.~Peres and B.~Vir\'ag, \emph{Zeros of Gaussian Analytic Functions and
Determinantal Point Processes}, University Lecture Series~51, American Mathematical Society, Providence, RI, 2009,
\doi{10.1090/ulect/051}.
\bibitem{JT98} B.~Jancovici and G.~T\'ellez, \emph{Two-dimensional Coulomb systems on a surface of constant negative
curvature}, J. Statist. Phys. \textbf{91} (1998), no.~5--6, 953--977, \doi{10.1023/A:1023079916489}.
\bibitem{Kri09} M.~Krishnapur, \emph{From random matrices to random analytic functions}, Ann. Probab. \textbf{37}
(2009), no.~1, 314--346, \doi{10.1214/08-AOP404}, \arxiv{0711.1378}.
\bibitem{Kal17} O.~Kallenberg, \emph{Random Measures, Theory and Applications}, Probability Theory and Stochastic
Modelling~77, Springer, Cham, 2017, \doi{10.1007/978-3-319-41598-7}.
\bibitem{PV05} Y.~Peres and B.~Vir\'ag, \emph{Zeros of the i.i.d.\ Gaussian power series: a conformally invariant
determinantal process}, Acta Math. \textbf{194} (2005), no.~1, 1--35, \doi{10.1007/BF02392515}.
\bibitem{MathBrainRepo} Y.~Shmalo, \emph{Math Brain}, public software repository, snapshot of 14 November 2026,
\url{https://github.com/yspennstate/math-brain-public/tree/8c9a6bd478c4ec4082347151b7ae16db50c917dd}.
\bibitem{MathBrainPaper} Y.~Shmalo, \emph{Math Brain: shared, auditable memory for multi-agent mathematical research},
manuscript dated 26 October 2026, Zenodo release 12 November 2026,
\url{https://github.com/yspennstate/math-brain-public/tree/main/papers/math-brain-2026-10-24},
\doi{10.5281/zenodo.24150837}.
\bibitem{BrainFramework} Y.~Shmalo, \emph{Brain Framework}, public software repository, snapshot of 13 November 2026,
\url{https://github.com/yspennstate/brain-framework-public/tree/5629e644cdd510173ae97e9aa92642dabd22dc57}.
\end{thebibliography}

\end{document}
